% Transcription — fonds Alexandre Grothendieck, shelfmark 119, batch 13 (pages 241-260).
% Established from the facsimile archives/batches/119/batch-13.pdf (a mirror of
% https://grothendieck.umontpellier.fr/119.pdf), per the skill
% transcribe-grothendieck.
%
% Pass: Opus 5.5 (claude-opus-5-5), 2026-10-03 — first pass, unchecked against the pages by a human.
%
% Folder 119 is the dossier for his CNRS application (1984). This batch holds
% no mathematics: it is his 1968 text on the university and a covering letter.
% Items (special rules for folder 119, issue #45):
%
%   pp. 241-250 — « Le maître-enseignant et le maître-chercheur dans
%     l'Université d'aujourd'hui », by A. Grothendieck, [June 1968]
%     (dated by the letter of pp. 251-252, which encloses it under this
%     title). His typescript, typed pages 1-10, with his manuscript
%     additions: « et de demain » in the title, the heading « A)
%     Aujourd'hui », and on p. 249 a recast of the closing paragraph
%     (« Je passerai en revue plus bas… », « Mon but dans la première
%     partie de cet exposé a été de mettre en évidence… »), plus a few
%     words in ink on pp. 247, 249, 250. TRANSCRIBED in full.
%   pp. 251-252 — his typed letter, signed by hand, to Hubert Beuve-Méry,
%     director of Le Monde, « Massy, le 4 Juin 1968 », offering the text
%     above for publication. His own letter, not printed anywhere.
%     TRANSCRIBED in full.
%   p. 253 — blank archivists' divider; no \page{}.
%   pp. 254-260 — the second state: the same text retyped, with the
%     manuscript additions of the first state taken into the type, typed
%     pages (1)-7, continuing in batch 14 to p. 270. TRANSCRIBED in full.
%
% The map given for this batch agrees with the pages.
%
% Differences seen between the two states over pp. 241-247 / 254-260:
%   - « et de demain » and « A. Aujourd'hui » typed in the second state;
%   - p. 257: « dans une certaine mesure l'insuffisance des crédits »
%     where p. 244 has « dans une large mesure »;
%   - p. 259: a footnote (*) on « bon enseignant » (« Du moins suivant la
%     notion traditionnelle… »), absent from p. 246;
%   - p. 259: « ni de temps » for « ni du temps », « se retrouver » for
%     « s'y retrouver »; p. 256: « l'éclat d'esprit » for « l'état
%     d'esprit » (a typing slip); « maîtres-assistants » hyphenated once;
%   - new typing slips: « tradionnellement », « tradionnels »,
%     « structure scolaires »; ink corrections of the first state (« de
%     futurs », « au manque », « pallier l' », « au contact de ») adopted.
%   - p. 254 keeps « comme l'enseignant », which p. 241 strikes to
%     « comme enseignant ».
%   The expansion proper (part B, « demain ») lies beyond this batch.
%
% Editorial conventions: words he underlined are set in \emph{}; typed
% quotation marks "…" are given as « … »; accents added in ink to the
% typescript and typist's overstrikes are silently adopted; letters or
% words struck in ink are \struck{}, ink insertions \add{}.
%
\documentclass[11pt,a4paper]{article}
\input{../preamble/grothendieck}

\folder{119}
\batch{13}
\pages{241}{260}
\dating{1968-1991}
\watermark{Édition de démonstration}
\foldertitle{Esquisse d’un programme [dossier constitué en vue d'une candidature au CNRS (1984)] : copies de tapuscrits et de tapuscrits annotés (1972-1974, 1978-1979, 1984, 1990-1991, s.d.), tirés à part (1971), notes et copies de note manuscrites (1986, s.d.), lettres (1968, 1972, 1991).}

\begin{document}

\section*{Le maître-enseignant et le maître-chercheur (juin 1968)}

\page{241}
\note{pagination dactylographiée de l'auteur : 1 à 10, p. 241 à 250}

\begin{quote}
Le maître-enseignant et le maître-chercheur dans l'Université d'aujourd'hui \add{et de demain}

par A. Grothendieck
\end{quote}

\subsection*{\add{A) Aujourd'hui.}}

1. Dans le système universitaire tel qu'il a existé jusqu'à nos jours,
depuis semble-t-il ses origines lointaines, le rôle du maître, d'une part
comme \struck{l'}enseignant chargé de transmettre des connaissances acquises, et
d'autre part comme créateur intellectuel contribuant à l'approfondissement,
à l'élargissement et au renouvellement de ces mêmes connaissances, était
confondu. En d'autres termes, une même personne était censée remplir ces
deux rôles, qui sont en effet étroitement apparentés, mais néanmoins
essentiellement distincts. N'était censé être qualifié pour enseigner la Science
que celui qui pouvait prétendre faire partie de ceux qui contribuaient à
l'édifier.

Ce système était effectivement compatible avec les besoins de la
société et de la science pendant fort longtemps. Les études universitaires
étaient réservées à un petit nombre, et les étudiants dans chaque discipline
poursuivaient généralement leurs études à un niveau théorique élevé, qui
les menait jusqu'aux frontières mouvantes entre le connu et le domaine de
la recherche. Il était naturel que l'enseignement universitaire soit assuré
par des maîtres dont chacun dominait à peu près complètement la science qu'il
enseignait, et dont chacun était lui-même un des artisans de cette science.
Vu le nombre limité d'étudiants, des effectifs modestes suffisaient pour
les enseigner, et le nombre des savants de valeur susceptibles d'assurer
cet enseignement était largement suffisant. Il n'était pas rare, même, que
des savants éminents soient obligés d'attendre de longues années la retraite
d'un de leurs collègues pour pouvoir occuper une position de titulaire de
chaire à la mesure de leur valeur.

2. Avec l'essor de la technologie dans ces dernières cinquante années, et
l'afflux de plus en plus considérable d'étudiants dans les Facultés, et
notamment dans les Facultés de Sciences, cette situation a radicalement

\page{242}
changé, non seulement en France, mais bien entendu partout dans le monde.
Tout d'abord, avec les \emph{motivations} de l'enseignement, qui deviennent de
plus en plus utilitaires, ses \emph{besoins} ont également changé. Parmi les
étudiants, de plus en plus nombreux sont ceux qui n'ont pas besoin d'une
connaissance théorique plus ou moins exhaustive, allant jusqu'aux limites
du connu, de telle ou telle science, la Mathématique disons ; mais plutôt
de telle ou telle partie de la Mathématique, qu'ils auront à utiliser
seulement à titre auxiliaire dans une variété de métiers techniques ou
scientifiques possibles : ingénieur, physicien, chimiste, biologiste, économiste…
Ce sont donc de\struck{s} futurs utilisateurs de Mathématiques, non de futurs
mathématiciens. Il leur est nécessaire de comprendre un formalisme mathématique
limité, en vue de pouvoir l'appliquer aux problèmes spécifiques de leurs
futurs métiers. Sous ces conditions, les questions de rigueur logique par
exemple (qui pour la Mathématique en tant que science pure sont et sans
doute resteront une exigence essentielle) deviennent accessoires, l'exigence
principale étant pour eux de savoir comprendre et surtout, appliquer
correctement, les algorithmes qu'ils seront amenés à utiliser.

Pour répondre à cette plus grande variété de besoins, il convient
donc de dispenser une plus grande variété d'enseignements,-- \emph{variété non
seulement par le nombre des théories mathématiques enseignées}, \emph{mais par l'optique
dans laquelle elles sont enseignées}, qui doit être fonction des besoins
véritables des étudiants, et ne pas rester soumise à des désiderata (de
rigueur, d'exhaustivité ou d'universalité) qui étaient justifiés il y a
cent ans et ne le sont plus aujourd'hui.

Parallèlement à cette différentiation des besoins des \emph{enseignés},
s'accentue nécessairement, qu'on le veuille ou non, la différentiation
entre les deux rôles dévolus au \emph{maître}. Sur le plan qualitatif, c'est
surtout vrai pour l'enseignement à assurer à de futurs utilisateurs de la
Mathématique comme plus haut, qui n'auront rien à gagner (et qui risquent
fort d'y perdre) que cet enseignement soit assuré par des savants éminents,
ou simplement des chercheurs de valeur. En règle générale et sauf de rares
exceptions, plus grande sera la valeur d'un chercheur, c'est-à-dire plus

\page{243}
fortement il sera aux prises avec les problèmes que lui pose la science qui
est en train de se faire, plus grand sera l'effort qu'il devra faire pour
s'arracher à ces problèmes, et la difficulté qu'il éprouvera pour se mettre
dans l'état d'esprit complètement différent que suppose un enseignement à
des étudiants qui ne s'intéressent nullement à cette science comme telle.
D'une part un tel enseignement a de fortes chances d'être faussé par les
exigences théoriques naturellement élevées du maître vis-à-vis de son propre
cours, lesquelles ne répondent nullement aux besoins de ses auditeurs,-- d'où
gaspillage de temps et d'énergie intellectuelle de ceux-ci. D'autre part, il
y a là un regrettable gaspillage d'énergie intellectuelle du maître également,
dont la fonction enseignante constitue dans le cas envisagé une entrave, et
non un stimulant ou un complément, à son activité créatrice.

Pour des raisons numériques également, on se trouve devant des
situations paradoxales et même absurdes tant qu'on continue à confondre
les rôles enseignant et créateur des maîtres à l'Université. Il est notoire,
pour toute personne bien informée de la question, que par suite de l'afflux
des étudiants et de la nécessité des enseignements divers à assurer, le
nombre des scientifiques qui sont capables d'un travail scientifique
original est de loin inférieur au nombre d'enseignants nécessaires à l'Université.
Ceci, joint au\struck{x} manque de crédits, a entraîné toute une série de graves
inconvénients qui sont allés en s'aggravant au fil des années :

a) Les anciens critères pour le recrutement des maîtres par leur
créativité scientifique étant théoriquement maintenus, un nombre
relativement petit et tout à fait insuffisant de maîtres se trouve être effectivement
recruté.

b) Pour pallier \struck{à} l'insuffisance croissante des effectifs de
maîtres de conférence et de professeurs, ceux-ci sont assistés dans leur
tâche par un nombre également croissant d'assistants et de maîtres assistants,
dont la position par rapport aux maîtres reste subalterne et la promotion
professionnelle problématique, par suite des critères mentionnés dans a).

c) Malgré l'appel aux forces enseignantes auxiliaires (mentionnées
dans b)) l'insuffisance des effectifs enseignants se fait durement sentir aux

\page{244}
étudiants comme aux enseignants, par des amphithéâtres surpeuplés, rendant
pratiquement impossible un contact pédagogique vivant entre enseignant et
enseigné.

Du point de vue des étudiants, c'est évidemment c) qui constitue
le vice principal, vice qui a sans doute été une des causes déterminantes
pour créer chez ceux-ci un état d'esprit qui a abouti à la salutaire crise
de l'Université que nous sommes en train de traverser. Alors que pour ce
dernier vice de fonctionnement, c'est dans une large mesure l'insuffisance
des crédits alloués qui semble être en cause, il semble que a) et b) doivent
être imputés surtout à la confusion entre le rôle du maître-enseignant et du
maître-chercheur. En fait, il se trouve que l'importance de la fonction
sociale du premier est généralement méconnue au profit du second, auquel
s'attache traditionnellement un prestige supérieur. Par d'anciennes habitudes
de pensée, le rôle du premier n'est considéré que comme un des attributs du
second. Contre l'enseignant universitaire qui n'est « \emph{que} » enseignant, il
y a dans le système actuel une prévention injustifiée et injuste, qui
vis-à-vis des maîtres assistants se traduit par une situation subalterne du
point de vue matériel comme du point de vue participation à la vie de
l'Université, et vis-à-vis des maîtres de conférence ou professeurs, à une
mésestime de la part de nombreux de leurs collègues, qui estiment (avec raison
parfois du point de vue des critères en vigueur jusqu'à maintenant) que
ceux-ci ont réussi à obtenir « au rabais » leur titre de docteur et leur
nomination à une chaire.

Le maintien de critères de recrutement et de valeur dépassés, et
l'état d'esprit qu'on vient de mentionner, créent d'ailleurs pour nombre
de scientifiques (et surtout de jeunes scientifiques) des problèmes
psychologiques et des problèmes de conscience tout à fait sérieux. Il y a en effet
le cas de scientifiques de valeur, ayant une culture étendue, un jugement
et un « goût » scientifiques très sûrs, des aptitudes pédagogiques excellentes,
qui sont animés d'un enthousiasme sincère pour leur science, mais qui
restent de longues années sans faire oeuvre vraiment originale, et qui n'y
parviennent parfois jamais. Alors que ces hommes seraient précieux pour

\page{245}
assurer un enseignement de très grande qualité, y compris à un niveau
élevé, dans une Université qui manque cruellement de telles forces, ils
se trouvent écartés de telles fonctions auxquelles tout les prédispose.
D'autres, parfaitement capables de faire des enseignants honorables, mais
n'ayant pas les dons voulus pour faire des travaux originaux, feront appel
à des « patrons » réputés particulièrement peu exigeants pour faire passer
une thèse, et avec l'appui de ceux-ci se verront nommés maîtres de conférence\struck{s},
alors que tel de leurs camarades, plus qualifié mais plus scrupuleux,
quittera le CNRS pour devenir simple assistant, ou cherchera un poste dans
l'enseignement secondaire. Par la suite, de tels maîtres de conférences,
qualifiés certes pour enseigner à l'Université, mais peu ou point pour juger
de questions concernant la formation ou le recrutement de mathématiciens
chercheurs, auront à ce sujet voix au chapitre au même titre que ceux de
leurs collègues ayant fait des contributions appréciables à la science, et
seront habilités pour faire passer des thèses (soi-disant de recherche)
tout comme ceux-ci. On voit que faute d'une définition des compétences qui
corresponde vraiment à la réalité, des hommes ayant une qualification
professionnelle sérieuse, qui pourrait sans usurpation aucune de leur part en
faire des membres utiles et à part entière de la communauté universitaire,
peuvent se trouver amenés, par la logique même du système dans lequel ils
sont placés et dont ils tirent parti à leur façon, dans une situation où
ils risquent de fausser la vie scientifique proprement dite dans leur
département, voire, dans des cas extrêmes, de l'étouffer.

3. Compte tenu de tout ce qui précède, il apparaît donc urgent de redéfinir,
les rôles de maître-enseignant et de maître-chercheur, pour les dissocier
tout en tenant compte en même temps de leurs interrelations étroites.

Le \emph{maître-chercheur} est un homme qui s'intéresse en premier lieu
aux problèmes que lui pose la science, plus qu'aux hommes qui l'apprennent
ou l'utilisent, plus même qu'aux hommes qui la font. L'enseignement souvent
pour lui est surtout un moyen, lui assurant une position sociale qui lui
permette de se consacrer à l'édification de la science indépendamment des
applications pratiques qui peuvent la rendre immédiatement rentable. Il est,

\page{246}
le plus souvent, un bon enseignant, et cela d'autant plus qu'il expose des
parties de la science qui sont d'un niveau théorique plus élevé et plus
proche des confins du connu. Il pourra être également bon enseignant au
niveau le plus élémentaire, traitant des bases même de la science, \emph{à condition}
qu'il s'adresse à un auditoire non pas de futurs utilisateurs sans plus,
mais à des étudiants qui se destinent à un approfondissement poussé de cette
science. Il n'est pas rare, sous cette réserve, que des savants de premier
ordre soient en même temps des enseignants émérites, et qu'ils aient pour
l'enseignement une passion qui augmente avec les années. Il arrive également,
mais c'est l'exception, qu'un chercheur de grande valeur soit un excécrable\note{sic}
enseignant : l'enseignement lui est une corvée qu'il accomplit à contre-coeur,
et il est autant à plaindre que ses étudiants, qui dès le début de l'année
académique perdent l'espoir de jamais pouvoir s'y retrouver dans les cours
qu'il dispense, touffus et laborieux quand ils ne sont obscurs et sibyllins.
Ils garderont sans doute le plus morne souvenir d'un homme dont ils n'ont
aucune raison de suspecter que c'est une des lumières de la science qu'il
leur enseigne si mal ! Personne, et lui-même moins que tout autre, ne
contestera que cet homme n'aurait jamais dû enseigner.

Le test le plus caractéristique pour distinguer le chercheur
mathématicien de toute autre espèce d'individu, c'est que lorsqu'il rencontre
un congénère, que ce soit autour d'une table, dans un métro, dans un congrès
ou n'importe où que ce soit, il se met à parler aussitôt non de politique,
ni du temps, ni de ses élèves, ni de ses collègues, ni de ses patrons (quand
il en a), mais de problèmes mathématiques. L'enseignement l'intéresse ou
même le passionne par les problèmes techniques d'exposition qu'il soulève,
variables à l'infini suivant le sujet enseigné et les connaissances supposées
acquises, qui mettent en jeu de mille façons imprévues son ingéniosité et
sa virtuosité, et non en fonction de la personnalité propre de ses divers
auditeurs. On peut rapporter à ce sujet une anecdote bien connue (dont il
n'est pas question de garantir l'authenticité !), selon laquelle le
mathématicien C.L. Siegel, connu également pour son amour de l'enseignement,
un jour que l'unique auditeur qui lui restait (il faut croire que c'était un

\page{247}
cours avancé) avait été empêché de venir, aurait fait son cours devant une
salle vide.

4. Le \emph{maître-enseignant} est un homme qui s'intéresse à la science autant ou
plus par les liens humains qu'elle implique, crée ou développe, que par
l'attrait de la découverte scientifique. Sur le plan de la fonction sociale,
c'est donc bien par le rôle d'enseignant qu'il pourra à présent le plus
naturellement rentabiliser sa vocation naturelle. Le plus souvent, il a des
dispositions pour la recherche, qui seules jusqu'à présent étaient censées
l'habiliter à suivre sa vocation principale. Mais souvent aussi, une fois sa
thèse terminée, qui lui aura demandé (à supposer que lui-même ou son
« patron de thèse » aient pris au sérieux les critères de valeur traditionnels
pour la thèse) un effort très considérable et unique dans sa vie, il
n'apporte plus guère d'autres contributions originales à sa science. Il arrive
même, faute d'aucun mécanisme de contrôle par ses étudiants ni par ses pairs,
qu'il perde tout contact avec la science vivante et qu'il s'installe
confortablement dans son statut de fonctionnaire, répétant au long des années
jusqu'à l'âge de la retraite, sans y changer une virgule, le même cours qu'il
aura mis au point dans sa jeunesse, quand il ne l'a pas hérité tel quel
de son vénérable prédécesseur. Mais nous sortons là de notre description du
maître-enseignant pour celle, trop connue hélas à tous ceux qui ont passé
par l'Université pour qu'il soit utile de la poursuivre, du \emph{maître-fossile}.
Ce dernier, véritable honte de notre métier, et fléau de l'enseignement
depuis l'école communale jusqu'à l'Université, n'est sans doute que
partiellement responsable de son propre rôle stérilisant, n'étant qu'un
sous-produit de structures scolaires et universitaires qui rendent possible le
foisonnement de lui et de ses pareils. Le vrai enseignant sait d'instinct
qu'il ne saura remplir valablement son rôle qu'à condition d'un renouvellement
constant dans les sujets enseignés, la façon de les présenter, voire même
dans les méthodes d'enseignement.

Pour un enseignant universitaire, un tel renouvellement n'est sans
doute possible qu'en restant au contact \struck{avec} \add{de} la science vivante. Ce n'est
qu'un tel contact qui gardera à son esprit cette ouverture et cette attitude
critique vis-à-vis de la science qu'il enseigne, qui caractérisent la « tête

\page{248}
bien faite » par opposition à la « tête bien pleine ». Cela ne suppose pas
nécessairement que le maître-enseignant soit au courant, à chaque moment,
de tous les plus récents progrès de sa spécialité, ce qui demanderait de sa
part un effort de mise au courant souvent incompatible avec ses tâches
d'enseignement. Mais les longues vacances universitaires lui permettent
tout au moins de se mettre au courant périodiquement des plus importants
parmi ces progrès, ou de contribuer, dans la mesure de ses moyens, à
l'avancement de sa science. Ces moyens personnels sont évidemment très
variables de l'un à l'autre. Au point de vue du maître-enseignant où nous
nous plaçons ici, l'étendue des dons pour la recherche devient d'ailleurs
un point secondaire, dans la mesure où sa participation à l'édification de
la science n'est pas conçue comme une fin en soi, mais comme un moyen parmi
d'autres pour garder vivant l'enseignement,-- étant entendu que cela suppose
d'abord que l'enseignant lui-même reste intellectuellement vivant,
c'est-à-dire intellectuellement actif.

Dans tout cela, la vocation principale reste l'enseignement, et
en définitive \emph{le sujet principal d'intérêt du maître-enseignant est l'étudiant
lui-même}. Il est clair qu'une telle vocation ne pourra s'épanouir pleinement
que sous des conditions qui permettent un contact vivant, constant et étroit
avec les étudiants. On sait combien nous en avons été loins\note{sic} tout au long de
ces dernières années, avec des amphithéâtres-forums où le
maître-gardien-de-la-science, seul devant la foule des aspirants-à-la-science, semble sorti
tout droit d'une hallucinante vision de Kafka.

5. Les développements qui précèdent démontrent sans doute suffisamment la
nécessité de distinguer entre le maître qui enseigne et le maître qui édifie
la science, et mettent en évidence quelques uns de leurs traits
caractéristiques. Il doit être bien entendu qu'il ne peut être question d'établir entre
ces deux espèces de vocations une hiérarchie quelconque. Il convient
simplement de réagir enfin à d'anciennes habitudes de pensée héritées du passé,
qui minimisaient le rôle du premier en le subordonnant au second. Il faut
donc réhabiliter la vocation pédagogique dans l'Université, en la mettant
à pied d'égalité avec la vocation de recherche, et en lui fournissant les

\page{249}
conditions favorables à son épanouissement qui ont jusqu'à présent fait
défaut.

En plus des maîtres eux-mêmes, les étudiants ont été évidemment les
principales victimes de la survie d'un système de formation et de
recrutement des maîtres basé presque exclusivement sur le critère de l'aptitude
à la recherche. Il est urgent de mettre à profit dans l'immédiat la force
vive libérée aujourd'hui par le mouvement de revendications étudiantes, et
les prises de conscience qu'elles ont précipité\add{es} chez les enseignants,
pour mettre au point des mécanismes assez souples pour distinguer et
coordonner tout à la fois les deux principales vocations de l'Université, et capables
de s'adapter au fur et à mesure aux besoins, changeants d'une période à l'autre
et d'une Faculté à l'autre. Ces questions sont l'objet dès à présent de
discussions entre étudiants et enseignants (assistants, maîtres assistants,
maîtres de conférences et professeurs), tout au moins au département de
Mathématiques de la Faculté des Sciences d'Orsay, et sans doute un peu
partout dans les universités en effervescence créatrice. Il faut que nous,
les principaux concernés, étudiants et maîtres enseignants et chercheurs,
fassions tout notre possible pour que des solutions d'ensemble soient
dégagées dans le plus proche avenir, et puissent être mises en application dès
la rentrée prochaine. \struck{Il n'est pas dans mon propos ici de passer}
\add{Je passerai} en revue \struck{les} \add{plus bas} diverses modifications en
profondeur des structures et des habitudes de
pensée dont la nécessité \add{me} semble dès à présent acquise. Certaines parmi les
plus importantes découlent logiquement de l'état de faits que je viens
d'essayer d'esquisser. Cet état de faits est bien familier sans doute à ceux
de mes collègues ayant assez d'expérience de l'enseignement ou de la
recherche, mais moins familier déjà, semble-t-il, aux assistants ou aux étudiants,
et a fortiori aux personnes étrangères à l'Université. Mon but dans \struck{l'}
\add{la première partie de cet} exposé \struck{qui précède n'a pas été de proposer
des \emph{solutions}, laissant ce soin à telles assemblées spontanées plus
autorisées que moi qui se réuniront pour en dégager en commun, mais de contribuer
à mettre} \add{a été de mettre} en évidence le plus clairement \struck{que je peux}
\add{possible} un des vices de base du système actuel, et \uncertain{d'}inciter tous
\note{le « d' » est surfrappé à la machine, peut-être sur un « à »}

\page{250}
ceux qui sont concernés à réfléchir à cette question pour contribuer à
l'élaboration de solutions novatrices. Les incertitudes du moment présent
sont plus lourdes de promesses que toutes les certitudes sur lesquelles se
reposait dame Université jusqu'à il y a peu. C'est de nous tous qui sommes
concernés, étudiants, enseignants et chercheurs, non des princes qui nous
gouvernent, qu'il dépend \struck{si} \add{que} ces promesses \struck{vont}
débouche\struck{r}\add{nt} sur cette rénovation de l'enseignement dont
personne (sujets ni princes) ne songe plus à contester la nécessité.

\note{fin du premier état ; la « partie B) » qu'annonce le titre ajouté à la main (« A) Aujourd'hui ») ne suit pas ici}

\section*{Lettre à Hubert Beuve-Méry, 4 juin 1968}

\page{251}
A. GROTHENDIECK

2, Avenue de Verrières

91 -- \emph{MASSY}

A Monsieur Hubert Beuve-Méry

Directeur du Monde

5, Rue des Italiens

\emph{PARIS} IX

Massy, le 4 Juin 1968

Cher Monsieur,

J'ai été invité par mes collègues du département de Mathématiques
de la Faculté des Sciences d'Orsay à participer aux travaux de la commission
chargée par ce département d'étudier les questions concernant la formation
et le recrutement des enseignants et chercheurs mathématiciens dans
l'Université. Au cours de plusieurs discussions entre étudiants, assistants ou
maîtres assistants, et professeurs, il est apparu que certains vices
fondamentaux du système actuel, dont à peu près tous les professeurs concernés
avaient une vision assez claire, étaient moins apparents pour les étudiants
et les assistants, qui sont pourtant directement concernés. Pour leur
permettre de s'associer en pleine connaissance de cause aux propositions et aux
décisions qui doivent être élaborés à ce sujet dans un prochain avenir, il
m'a semblé qu'il était nécessaire d'informer sur ces problèmes un public
plus large que celui formé des quelques delégués d'étudiants et d'assistants
qui sont amenés à prendre part à ces discussions, qui devrait comprendre pour
le moins l'ensemble des étudiants et enseignants en mathématiques, ou même
en sciences en général,-- les problèmes soulevés étant sans doute les mêmes
(quoique probablement moins aigus) dans les autres sciences. C'est dans ce
but que je vous envoie ci-joint un exposé intitulé « Le maître-enseignant et
le maître-chercheur dans l'Université d'aujourd'hui », en vous suggérant la
publication de cet exposé, sous forme de un ou deux articles de fond, dans
un avenir rapproché. Je vous serais obligé de me faire connaître, dès qu'il
vous sera possible, s'il vous sera possible de publier cet exposé dans une

\page{252}
de vos prochaines éditions, pour me permettre dans le cas contraire de
prévoir d'autres modes de diffusion rapides de cet exposé, que je tiens
à mettre à la disposition du public concerné dans les plus brefs délais.
A cause de la grève des postes, je vous prie de me faire connaître votre
réponse par téléphone à mon numéro personnel, 920 13 34.

En vous remerciant d'avance, je vous prie d'agréer, Monsieur,
l'expression de ma plus parfaite considération.

\add{A Grothendieck}\note{signature manuscrite}

A. Grothendieck

Professeur à l'Institut des Hautes Etudes Scientifiques

91 -- Bures sur Yvette

Tél : 928 48 53

\section*{Second état, retapé et augmenté}

\page{254}
\note{en haut à droite, au crayon, d'une main qu'on ne peut attribuer : « 4, 7, 11, \ill{} 7 », coupé par une déchirure du feuillet ; cette page n'est pas numérotée à la machine}

\begin{quote}
Le maître-enseignant et le maître-chercheur dans l'Université d'aujourd'hui et de demain

par A. Grothendieck
\end{quote}

\subsection*{A. \emph{Aujourd'hui}}

1. Dans le système universitaire tel qu'il a existé jusqu'à nos jours, depuis
semble-t-il ses origines lointaines, le rôle du maître, d'une part comme
l'enseignant chargé de transmettre des connaissances acquises, et d'autre part
comme créateur intellectuel contribuant à l'approfondissement, à
l'élargissement et au renouvellement de ces mêmes connaissances, était confondu. En
d'autres termes, une même personne était censée remplir ces deux rôles, qui
sont en effet étroitement apparentés, mais néanmoins essentiellement distincts.
N'était censé être qualifié pour enseigner la Science que celui qui pouvait
prétendre faire partie de ceux qui contribuaient à l'édifier.

Ce système était effectivement compatible avec les besoins de la
société et de la science pendant fort longtemps. Les études universitaires
étaient réservées à un petit nombre, et les étudiants dans chaque discipline
poursuivaient généralement leurs études à un niveau théorique élevé, qui les
menait jusqu'aux frontières mouvantes entre le connu et le domaine de la
recherche. Il était naturel que l'enseignement universitaire soit assuré par
des maîtres dont chacun dominait à peu près complètement la science qu'il
enseignait, et dont chacun était lui-même un des artisans de cette science.
Vu le nombre limité d'étudiants, des effectifs modestes suffisaient pour les
enseigner, et le nombre des savants de valeur susceptibles d'assurer cet
enseignement était largement suffisant. Il n'était pas rare, même, que des
savants éminents soient obligés d'attendre de longues années la retraite d'un
de leurs collègues pour pouvoir occuper une position de titulaire de chaire
à la mesure de leur valeur.

2. Avec l'essor de la technologie dans ces dernières cinquante années, et
l'afflux de plus en plus considérable d'étudiants dans les Facultés, et
notamment dans les Facultés de Sciences, cette situation a radicalement changé, non

\page{255}
seulement en France, mais bien entendu partout dans le monde. Tout d'abord,
avec les \emph{motivations} de l'enseignement, qui deviennent de plus en plus
utilitaires, ses \emph{besoins} ont également changé. Parmi les étudiants, de plus en
plus nombreux sont ceux qui n'ont pas besoin d'une connaissance théorique plus
ou moins exhaustive, allant jusqu'aux limites du connu, de telle ou telle
science, la Mathématique disons ; mais plutôt de telle ou telle partie de la
Mathématique, qu'ils auront à utiliser seulement à titre auxiliaire dans une
variété de métiers techniques ou scientifiques possibles : ingénieur,
physicien, chimiste, biologiste, économiste… Ce sont donc de futurs utilisateurs
de Mathématiques, non de futurs mathématiciens. Il leur est nécessaire de
comprendre un formalisme mathématique limité, en vue de pouvoir l'appliquer
aux problèmes spécifiques de leurs futurs métiers. Sous ces conditions, les
questions de rigueur logique par exemple (qui pour la Mathématique en tant
que science pure sont et sans doute resteront une exigence essentielle)
deviennent accessoires, l'exigence principale étant pour eux de savoir comprendre
et surtout, appliquer correctement, les algorithmes qu'ils seront amenés à
utiliser.

Pour répondre à cette plus grande variété de besoins, il convient
donc de dispenser une plus grande variété d'enseignements,-- \emph{variété non
seulement par le nombre des théories mathématiques enseignées}, \emph{mais par l'optique
dans laquelle elles sont enseignées}, qui doit être fonction des besoins
véritables des étudiants, et ne pas rester soumise à des desiderata (de
rigueur, d'exhaustivité ou d'universalité) qui étaient justifiés il y a cent
ans et ne le sont plus aujourd'hui.

Parallèlement à cette différentiation des besoins des \emph{enseignés},
s'accentue nécessairement, qu'on le veuille ou non, la différentiation entre
les deux rôles dévolus au \emph{maître}. Sur le plan qualitatif, c'est surtout
vrai pour l'enseignement à assurer à de futurs utilisateurs de la
Mathématique comme plus haut, qui n'auront rien à gagner (et qui risquent fort d'y
perdre) que cet enseignement soit assuré par des savants éminents, ou
simplement des chercheurs de valeur. En règle générale et sauf de rares
exceptions, plus grande sera la valeur d'un chercheur, c'est-à-dire plus fortement

\page{256}
il sera aux prises avec les problèmes que lui pose la science qui est en train
de se faire, plus grand sera l'effort qu'il devra faire pour s'arracher à ces
problèmes, et la difficulté qu'il éprouvera pour se mettre dans l'éclat\note{sic, pour « l'état » du premier état, p. 243}
d'esprit complètement différent que suppose un enseignement à des étudiants qui
ne s'intéressent nullement à cette science comme telle. D'une part un tel
enseignement a de fortes chances d'être faussé par les exigences théoriques
naturellement élevées du maître vis-à-vis de son propre cours, lesquelles
ne repondent nullement aux besoins de ses auditeurs,-- d'où gaspillage de
temps et d'énergie intellectuelle de ceux-ci. D'autre part, il y a là un
regrettable gaspillage d'énergie intellectuelle du maître également, dont
la fonction enseignante constitue dans le cas envisagé une entrave, et non
un stimulant ou un complément, à son activité créatrice.

Pour des raisons numériques également, on se trouve devant des
situations paradoxales et même absurdes tant qu'on continue à confondre les
rôles enseignant et créateur des maîtres à l'Université. Il est notoire,
pour toute personne bien informée de la question, que par suite de l'afflux
des étudiants et de la nécessité des enseignements divers à assurer, le
nombre des scientifiques qui sont capables d'un travail scientifique original
est de loin inférieur au nombre d'enseignants nécessaires à l'Université.
Ceci, joint au manque de crédits, a entraîné toute une série de graves
inconvénients qui sont allés en s'aggravant au fil des années :

a) Les anciens critères pour le recrutement des maîtres par leur
créativité scientifique étant théoriquement maintenus, un nombre
relativement petit et tout à fait insuffisant de maîtres se trouve être
effectivement recruté.

b) Pour pallier l'insuffisance croissante des effectifs de maîtres
de conférence et de professeurs, ceux-ci sont assistés dans leur tâche par
un nombre également croissant d'assistants et de maîtres-assistants, dont la
position par rapport aux maîtres reste subalterne et la promotion
professionnelle problématique, par suite des critères mentionnés dans a).

c) Malgré l'appel aux forces enseignantes auxiliaires (mentionnées
dans b)) l'insuffisance des effectifs enseignants se fait durement sentir aux

\page{257}
étudiants comme aux enseignants, par des amphithéâtres surpeuplés, rendant
pratiquement impossible un contact pédagogique vivant entre enseignant et
enseigné.

Du point de vue des étudiants, c'est évidemment c) qui constitue le
vice principal, vice qui a sans doute été une des causes déterminantes pour
créer chez ceux-ci un état d'esprit qui a abouti à la salutaire crise de
l'Université que nous sommes en train de traverser. Alors que pour ce dernier
vice de fonctionnement, c'est dans une certaine mesure l'insuffisance des crédits
alloués qui semble être en cause, il semble que a) et b) doivent être imputés
surtout à la confusion entre le rôle du maître-enseignant et du maître-chercheur.
En fait, il se trouve que l'importance de la fonction sociale du premier est
généralement méconnue au profit du second, auquel s'attache tradionnellement\note{sic}
un prestige supérieur. Par d'anciennes habitudes de pensée, le rôle du premier
n'est considéré que comme un des attributs du second. Contre l'enseignant
universitaire qui n'est « \emph{que} » enseignant, il y a dans le système actuel une
prévention injustifiée et injuste, qui vis-à-vis des maîtres assistants se
traduit par une situation subalterne du point de vue matériel comme du point
de vue participation à la vie de l'Université, et vis-à-vis des maîtres de
conférence ou professeurs, à une mésestime de la part de nombreux de leurs
collègues, qui estiment (avec raison parfois du point de vue des critères
en vigueur jusqu'à maintenant) que ceux-ci ont réussi à obtenir « au rabais »
leur titre de docteur et leur nomination à une chaire.

Le maintien de critères de recrutement et de valeur dépassés, et
l'état d'esprit qu'on vient de mentionner, créent d'ailleurs pour nombre de
scientifiques (et surtout de jeunes scientifiques) des problèmes psychologiques
et des problèmes de conscience tout à fait sérieux. Il y a en effet le cas
de scientifiques de valeur, ayant une culture étendue, un jugement et un
« goût » scientifiques très sûrs, des aptitudes pédagogiques excellentes, qui
sont animés d'un enthousiasme sincère pour leur science, mais qui restent
de longues années sans faire oeuvre vraiment originale, et qui n'y parviennent
parfois jamais. Alors que ces hommes seraient précieux pour assurer un

\page{258}
enseignement de très grande qualité, y compris à un niveau élevé, dans une
Université qui manque cruellement de telles forces, ils se trouvent écartés de telles
fonctions auxquelles tout les prédispose. D'autres, parfaitement capables de
faire des enseignants honorables, mais n'ayant pas les dons voulus pour faire
des travaux originaux, feront appel à des « patrons » réputés particulièrement
peu exigeants pour faire passer une thèse, et avec l'appui de ceux-ci se
verront nommés maîtres de conférence\struck{s}, alors que tel de leurs camarades, plus
qualifié mais plus scrupuleux, quittera le CNRS pour devenir simple assistant,
ou cherchera un poste dans l'enseignement secondaire. Par la suite, de tels
maîtres de conférences, qualifiés certes pour enseigner à l'Université, mais
peu ou point pour juger de questions concernant la formation ou le recrutement
de mathématiciens chercheurs, auront à ce sujet voix au chapitre au même
titre que ceux de leurs collègues ayant fait des contributions appréciables
à la science, et seront habilités pour faire passer des thèses (soi-disant de
recherche) tout comme ceux-ci. On voit que faute d'une définition des
compétences qui corresponde vraiment à la réalité, des hommes ayant une
qualification professionnelle sérieuse, qui pourrait sans usurpation aucune de leur
part en faire des membres utiles et à part entière de la communauté
universitaire, peuvent se trouver amenés, par la logique même du système dans
lequel ils sont placés et dont ils tirent parti à leur façon, dans une
situation où ils risquent de fausser la vie scientifique proprement dite
dans leur département, voire, dans des cas extrêmes, de l'étouffer.

3. Compte tenu de tout ce qui précède, il apparaît donc urgent de redéfinir,
les rôles de maître-enseignant et de maître-chercheur, pour les dissocier
tout en tenant compte en même temps de leurs interrelations étroites.

Le \emph{maître-chercheur} est un homme qui s'intéresse en premier lieu
aux problèmes que lui pose la science, plus qu'aux hommes qui l'apprennent
ou l'utilisent, plus même qu'aux hommes qui la font. L'enseignement souvent
pour lui est surtout un moyen, lui assurant une position sociale qui lui
permette de se consacrer à l'édification de la science indépendamment des
applications pratiques qui peuvent la rendre immédiatement rentable. Il est,

\page{259}
le plus souvent, un bon enseignant (*), et cela d'autant plus qu'il expose des
parties de la science qui sont d'un niveau théorique plus élevé et plus
proche des confins du connu. Il pourra être également bon enseignant au niveau
le plus élémentaire, traitant des bases même de la science, \emph{à condition}
qu'il s'adresse à un auditoire non pas de futurs utilisateurs sans plus,
mais à des étudiants qui se destinent à un approfondissement poussé de cette
science. Il n'est pas rare, sous cette réserve, que des savants de premier
ordre soient en même temps des enseignants émérites, et qu'ils aient pour
l'enseignement une passion qui augmente avec les années. Il arrive également,
mais c'est l'exception, qu'un chercheur de grande valeur soit un excécrable\note{sic}
enseignant : l'enseignement lui est une corvée qu'il accomplit à contre-coeur,
et il est autant à plaindre que ses étudiants, qui dès le début de l'année
académique perdent l'espoir de jamais pouvoir se retrouver dans les cours
qu'il dispense, touffus et laborieux quand ils ne sont obscurs et sibyllins.
Ils garderont sans doute le plus morne souvenir d'un homme dont ils n'ont
aucune raison de suspecter que c'est une des lumières de la science qu'il
leur enseigne si mal ! Personne, et lui-même moins que tout autre, ne
contestera que cet homme n'aurait jamais dû enseigner.

Le test le plus caractéristique pour distinguer le chercheur
mathématicien de toute autre espèce d'individu, c'est que lorsqu'il rencontre
un congénère, que ce soit autour d'une table, dans un métro, dans un congrès
ou n'importe où que ce soit, il se met à parler aussitôt non de politique,
ni de temps, ni de ses élèves, ni de ses collègues, ni de ses patrons (quand
il en a), mais de problèmes mathématiques. L'enseignement l'intéresse ou
même le passionne par les problèmes techniques d'exposition qu'il soulève,
variables à l'infini suivant le sujet enseigné et les connaissances supposées
acquises, qui mettent en jeu de mille façons imprévues son ingéniosité et
sa virtuosité, et non en fonction de la personnalité propre de ses divers
auditeurs. On peut rapporter à ce sujet une anecdote bien connue (dont il
n'est pas question de garantir l'authenticité !), selon laquelle le
mathématicien C.L. Siegel, connu également pour son amour de l'enseignement,
un jour que l'unique auditeur qui lui restait (il faut croire que c'était un

(*) Du moins suivant la notion traditionnelle du « bon enseignant » ; en fait, il
y a des étudiants qui trouvent tel assistant meilleur enseignant que son
patron professeur, réputé bon enseignant parmi ses pairs !

\page{260}
cours avancé) avait été empêché de venir, aurait fait son cours devant une
salle vide.

4. Le \emph{maître-enseignant} est un homme qui s'intéresse à la science autant ou
plus par les liens humains qu'elle implique, crée ou développe, que par
l'attrait de la découverte scientifique. Sur le plan de la fonction sociale,
c'est donc bien par le rôle d'enseignant qu'il pourra à présent le plus
naturellement rentabiliser sa vocation naturelle. Le plus souvent, il a des
dispositions pour la recherche, qui seules jusqu'à présent étaient censées
l'habiliter à suivre sa vocation principale. Mais souvent aussi, une fois sa
thèse terminée, qui lui aura demandé (à supposer que lui-même ou son
« patron de thèse » aient pris au sérieux les critères de valeur tradionnels\note{sic}
pour la thèse) un effort très considérable et unique dans sa vie, il
n'apporte plus guère d'autres contributions originales à sa science. Il arrive
même, faute d'aucun mécanisme de contrôle par ses étudiants ni par ses pairs,
qu'il perde tout contact avec la science vivante et qu'il s'installe
confortablement dans son statut de fonctionnaire, répétant au long des années
jusqu'à l'âge de la retraite, sans y changer une virgule, le même cours qu'il
aura mis au point dans sa jeunesse, quand il ne l'a pas hérité tel quel
de son vénérable prédécesseur. Mais nous sortons là de notre description du
maître-enseignant pour celle, trop connue hélas à tous ceux qui ont passé
par l'Université pour qu'il soit utile de la poursuivre, du \emph{maître-fossile}.
Ce dernier, véritable honte de notre métier, et fléau de l'enseignement depuis
l'école communale jusqu'à l'Université, n'est sans doute que partiellement
responsable de son propre rôle stérilisant, n'étant qu'un sous-produit de
structure\note{sic} scolaires et universitaires qui rendent possible le foisonnement de
lui et de ses pareils. Le vrai enseignant sait d'instinct qu'il ne saura
remplir valablement son rôle qu'à condition d'un renouvellement constant dans
les sujets enseignés, la façon de les présenter, voire même dans les méthodes
d'enseignement.

Pour un enseignant universitaire, un tel renouvellement n'est sans
doute possible qu'en restant au contact de la science vivante. Ce n'est qu'un
tel contact qui gardera à son esprit cette ouverture et cette attitude
critique vis-à-vis de la science qu'il enseigne, qui caractérisent la « tête
\note{la phrase et le texte se poursuivent au lot suivant (p. 261 et suiv., jusqu'à la p. 270)}

\end{document}
